\subsection{محدودیت‌ها و تحلیل شکست}
اگرچه ماژول در مدل‌های انکودر و دیکودر استاندارد به وضعیت \lr{PASS} خالص دست یافته است، اما نکات تحلیلی زیر در کاربردهای وسیع‌تر شایان توجه است:
\begin{itemize}
    \item \textbf{کاهش توان عملیاتی در مدل‌های بزرگ‌تر:} توان پردازشی بر روی مدل \lr{RoBERTa} به \lr{4,678.15 tok/s} و بر روی \lr{GPT-2} به \lr{5,372.02 tok/s} کاهش یافته است. علت این امر، افزونگی محاسباتی ناشی از اجرای همزمان دو مسیر توجه (محلی و سراسری) به صورت متوالی و عملیات کِرنلی در پایتورچ بدون تلفیق عملگر سخت‌افزاری (\lr{Kernel Fusion}) در سطح \lr{CUDA} است.
    \item \textbf{حساسیت در مدل‌های زبان علّی فوق حجیم (\lr{Qwen}) و \lr{T5}:} در مدل \lr{Qwen} اگرچه با دقت \lr{53.79\%} وضعیت \lr{PASS} کسب گردیده، اما مصرف حافظه به \lr{1,195.99 MB} رسیده است که ناشی از سربار بافر تانسورهای پنهان در لایه‌های عمیق دکودر است. همچنین به دلیل عدم تطابق سازوکار توجه متقاطع انکودر-دیکودر با پروجکشن اختصاصی، تست روی مدل \lr{T5} با وضعیت \lr{SKIP\_T5\_INCOMPATIBLE} برچسب‌گذاری شده است.
\end{itemize}
